% GENERATED FILE. Edit canonical lessons, then run npm run generate:books.
% canonical-source-hash: fnv1a64:40a04eaecd9ff8d8
% canonical-lessons: TE-C112-sanci, TE-C112-kalam, TE-C112-kagitam, TE-C112-glasu, TE-C112-cemca

\chapter{Bag, Pen, Paper, Drinking glass, Spoon}
\label{ch:te-bag-pen-paper-drinking-glass-spoon}
\hlchaptermodality{112}

\begin{chapteropening}
\textbf{By the end of this chapter:} \emph{I can say a bag, a pen, paper, a drinking glass and a spoon.}

Complete the last lesson of chapter 112: i can say a bag, a pen, paper, a drinking glass and a spoon.
\end{chapteropening}

\section[sañci]{\te{సంచి} (sañci) — a bag}

\label{lesson:TE-C112-sanci}



\begin{quote}
Before the new one: say the Telugu for glasses, then the Telugu for a clock.
\end{quote}

\subsection*{You'll want to know: \te{సంచి}}
\textbf{\te{సంచి}} — \emph{sañci} — \textquotedblleft{}a bag\textquotedblright{}.

\begin{grammarlens}[title={What you've built}]
One more word for this chapter.
\end{grammarlens}

\subsection*{Your turn}
\begin{itemize}
\raggedright
  \item \emph{Say it:} \emph{sañci}
  \item \emph{Say it:} \emph{sañci}, once more
  \item \emph{Say it:} say \emph{sañci}
  \item \emph{Recall:} say the Telugu for glasses, then the Telugu for a clock, then say \emph{sañci} again
\end{itemize}

\begin{tcolorbox}[breakable,colback=teal!4,colframe=teal!35!black,title={Before you move on}]
What does \te{సంచి} mean? (\textquotedblleft{}A bag\textquotedblright{}.) Say it once more.
\end{tcolorbox}
\section[kalaṁ]{\te{కలం} (kalaṁ) — a pen}

\label{lesson:TE-C112-kalam}



\begin{quote}
Before the new one: say the Telugu for a clock, then the Telugu for a bag.
\end{quote}

\subsection*{You'll want to know: \te{కలం}}
\textbf{\te{కలం}} — \emph{kalaṁ} — \textquotedblleft{}a pen\textquotedblright{}.

\begin{grammarlens}[title={What you've built}]
One more word for this chapter.
\end{grammarlens}

\subsection*{Your turn}
\begin{itemize}
\raggedright
  \item \emph{Say it:} \emph{kalaṁ}
  \item \emph{Say it:} \emph{kalaṁ}, once more
  \item \emph{Say it:} say \emph{kalaṁ}
  \item \emph{Recall:} say the Telugu for a clock, then the Telugu for a bag, then say \emph{kalaṁ} again
\end{itemize}

\begin{tcolorbox}[breakable,colback=teal!4,colframe=teal!35!black,title={Before you move on}]
What does \te{కలం} mean? (\textquotedblleft{}A pen\textquotedblright{}.) Say it once more.
\end{tcolorbox}
\section[kāgitaṁ]{\te{కాగితం} (kāgitaṁ) — paper}

\label{lesson:TE-C112-kagitam}



\begin{quote}
Before the new one: say the Telugu for a bag, then the Telugu for a pen.

Read the lines aloud as one run: \textbf{\te{ఇది} \te{పుస్తకం}. \te{కుర్చీ} \te{ఇక్కడ} \te{ఉంది}.}
\end{quote}

\subsection*{You'll want to know: \te{కాగితం}}
\textbf{\te{కాగితం}} — \emph{kāgitaṁ} — \textquotedblleft{}paper\textquotedblright{}.

\begin{grammarlens}[title={What you've built}]
One more word for this chapter.
\end{grammarlens}

\subsection*{Your turn}
\begin{itemize}
\raggedright
  \item \emph{Say it:} \emph{kāgitaṁ}
  \item \emph{Say it:} \emph{kāgitaṁ}, once more
  \item \emph{Say it:} say \emph{kāgitaṁ}
  \item \emph{Recall:} say the Telugu for a bag, then the Telugu for a pen, then say \emph{kāgitaṁ} again
\end{itemize}

\begin{tcolorbox}[breakable,colback=teal!4,colframe=teal!35!black,title={Before you move on}]
What does \te{కాగితం} mean? (\textquotedblleft{}Paper\textquotedblright{}.) Say it once more.
\end{tcolorbox}
\section[glāsu]{\te{గ్లాసు} (glāsu) — a drinking glass}

\label{lesson:TE-C112-glasu}



\begin{quote}
Before the new one: say the Telugu for a pen, then the Telugu for paper.
\end{quote}

\subsection*{You'll want to know: \te{గ్లాసు}}
\textbf{\te{గ్లాసు}} — \emph{glāsu} — \textquotedblleft{}a drinking glass\textquotedblright{}.

\begin{grammarlens}[title={What you've built}]
One more word for this chapter.
\end{grammarlens}

\subsection*{Your turn}
\begin{itemize}
\raggedright
  \item \emph{Say it:} \emph{glāsu}
  \item \emph{Say it:} \emph{glāsu}, once more
  \item \emph{Say it:} say \emph{glāsu}
  \item \emph{Recall:} say the Telugu for a pen, then the Telugu for paper, then say \emph{glāsu} again
\end{itemize}

\begin{tcolorbox}[breakable,colback=teal!4,colframe=teal!35!black,title={Before you move on}]
What does \te{గ్లాసు} mean? (\textquotedblleft{}A drinking glass\textquotedblright{}.) Say it once more.
\end{tcolorbox}
\section[cemcā]{\te{చెంచా} (cemcā) — a spoon}

\label{lesson:TE-C112-cemca}



\begin{quote}
Before the new one: say the Telugu for paper, then the Telugu for a drinking glass.
\end{quote}

\subsection*{You'll want to know: \te{చెంచా}}
\textbf{\te{చెంచా}} — \emph{cemcā} — \textquotedblleft{}a spoon\textquotedblright{}.

\begin{grammarlens}[title={What you've built}]
One more word for this chapter.
\end{grammarlens}

\subsection*{Your turn}
\begin{itemize}
\raggedright
  \item \emph{Say it:} \emph{cemcā}
  \item \emph{Say it:} \emph{cemcā}, once more
  \item \emph{Say it:} say \emph{cemcā}
  \item \emph{Recall:} say the Telugu for paper, then the Telugu for a drinking glass, then say \emph{cemcā} again
\end{itemize}

\begin{tcolorbox}[breakable,colback=teal!4,colframe=teal!35!black,title={Before you move on}]
What does \te{చెంచా} mean? (\textquotedblleft{}A spoon\textquotedblright{}.) Say it once more.
\end{tcolorbox}
